% Part: lambda-calculus
% Chapter: church-rosser
% Section: parallel-beta-eta-reduction

\documentclass[../../../include/open-logic-section]{subfiles}

\begin{document}

\olfileid[hi]{lam}{cr}{pbe}

\olsection{समांतर $\beta\eta$-अपचयन}

इस खंड में हम समांतर $\beta\eta$-अपचयन के लिए चर्च--रॉसर गुण सिद्ध करते
हैं; यह $\beta\eta$-अपचयन के अनुरूप समांतर अपचयन की धारणा है।

\begin{defn}[समांतर $\beta\eta$-अपचयन, $\beredpar$] \ollabel{defn:beredpar}
  पदों पर \emph{समांतर $\beta\eta$-अपचयन} ($\beredpar$) आगमनतः इस प्रकार
  परिभाषित है:
  \begin{enumerate}
    \item \ollabel{defn:beredpar1} $x \beredpar x$.
    \item \ollabel{defn:beredpar2} यदि $N \xrightarrow{\beta} N'$, तो $\lambd[x][N] \beredpar
      \lambd[x][N']$।
    \item \ollabel{defn:beredpar3} यदि $P \beredpar P'$ और $Q \beredpar Q'$, तो $PQ \beredpar
      P'Q'$।
    \item \ollabel{defn:beredpar4} यदि $N \beredpar N'$ और $Q \beredpar Q'$, तो
      $(\lambd[x][N])Q \beredpar \Subst{N'}{Q'}{x}$.
    \item \ollabel{defn:beredpar5} यदि $N \beredpar N'$, तो $\lambd[x][Nx]
      \beredpar N'$, बशर्ते $x \notin FV(N)$।
  \end{enumerate}
\end{defn}

\begin{thm}\ollabel{thm:refl}
  $M \beredpar M$.
\end{thm}

\begin{proof}
  अभ्यास।
\end{proof}

\begin{prob}
  \olref[lam][cr][pbe]{thm:refl} सिद्ध करें।
\end{prob}

\begin{defn}[$\beta\eta$-पूर्ण विकास]\ollabel{defn:becd}
  ~$M$ का \emph{$\beta\eta$-पूर्ण विकास}~$\becd{M}$ इस प्रकार परिभाषित है:
  \begin{align}
    \becd{x} &= x \ollabel{defn:becd1} \\
    \becd{(\lambd[x][N])} &= \lambd[x][\becd{N}] \ollabel{defn:becd2}\\
    \becd{(PQ)} &= \becd{P}\becd{Q} && \text{यदि $P$ कोई $\lambd$-अमूर्त नहीं है}
    \ollabel{defn:becd3} \\
    \becd{((\lambd[x][N])Q)} &= \Subst{\becd{N}}{\becd{Q}}{x}
                             \ollabel{defn:becd4} \\
    \becd{(\lambd[x][Nx])} &= \becd{N} \ollabel{defn:becd5} & \text{यदि $x
                                                      \notin FV(N)$}
  \end{align}
\end{defn}

\begin{lem}\ollabel{lem:comp}
  यदि $M \beredpar M'$ और $R \beredpar R'$, तो $\Subst{M}{R}{y}
  \beredpar \Subst{M'}{R'}{y}$।
\end{lem}

\begin{proof}
  $M \beredpar M'$ के !!{derivation} पर आगमन द्वारा।

  पहली चार स्थितियाँ ठीक \olref[pb]{lem:comp} वाली हैं। यदि अंतिम नियम
  \olref{defn:beredpar5} है, तो किन्हीं $x$ और~$N'$ के लिए $M$,
  ~$\lambd[x][Nx]$ और $M'$, ~$N'$ है, जहाँ $x \notin FV(N)$ और $N
  \beredpar N'$। हमें दिखाना है कि $\Subst{(\lambd[x][Nx])}{R}{y}
  \beredpar \Subst{N'}{R'}{y}$, अर्थात् $\lambd[x][\Subst{N}{R}{y} x]
  \beredpar \Subst{N'}{R'}{y}$। यह
  \olref{defn:beredpar}\olref{defn:beredpar5} और आगमन-परिकल्पना से मिलता
  है।
\end{proof}

\begin{lem}\ollabel{lem:cont}
  यदि $M \beredpar M'$, तो $M' \beredpar \becd{M}$।
\end{lem}

\begin{proof}
  $M \beredpar M'$ के !!{derivation} पर आगमन द्वारा।

  पहली चार स्थितियाँ \olref[pb]{lem:cont} वाली हैं। यदि अंतिम नियम
  \olref{defn:beredpar5} है, तो किन्हीं $x$, $N$, $N'$ के लिए $M$,
  ~$\lambd[x][Nx]$ और $M'$, ~$N'$ है, जहाँ $x \notin FV(N)$ और $N
  \beredpar N'$। हमें दिखाना है कि $N' \beredpar
  \becd{(\lambd[x][Nx])}$, अर्थात् $N' \beredpar \becd{N}$; यह
  आगमन-परिकल्पना से तात्कालिक है।
\end{proof}

\begin{thm}\ollabel{thm:cr}
  $\beredpar$ में चर्च--रॉसर गुण है।
\end{thm}

\begin{proof}
  \olref{lem:cont} से तात्कालिक।
\end{proof}

\end{document}
